\section{Introduction}

Fungi are prolific producers of secondary metabolites -- polyketides, non-ribosomal peptides,
terpenoids and hybrids thereof -- with roles spanning virulence, chemical ecology, and pharmacology,
from aflatoxins and other mycotoxins to clinically important natural products. The genes responsible
for a given metabolite are typically physically clustered in the genome as a biosynthetic gene
cluster (BGC), and the Minimum Information about a Biosynthetic Gene cluster (MIBiG) standard and
its accompanying database have become the community reference point for recording which BGCs have
been characterised and what they are reported to produce \citep{terlouw2023mibig}.

A MIBiG accession, however, is deliberately a minimal record: cluster boundaries, a product name,
and a small set of controlled fields. The primary evidence behind that record -- which individual
gene was actually disrupted and shown to be required, which metabolite congener was directly
detected versus inferred, under what culture conditions, with what caveats, and whether independent
studies agree -- lives only in the cited literature, in a different format for every paper, and is
not itself queryable. This creates a substantial gap between what is nominally known about a fungal
BGC and what can be programmatically retrieved, compared across BGCs, or checked for internal
disagreement.

\todo{Introduction -- add 2--4 discipline-specific citations here once verified (e.g. antiSMASH/BGC
genome-mining tools, prior fungal-BGC or natural-product knowledge-graph resources, and any of the
609 dossiers' own source papers that motivate specific examples). None are included yet because
fabricating bibliographic details from memory is deliberately avoided; see
\texttt{manuscript/latex/references.bib} for a note on this.}

The FAIR principles -- that data and the tools built on them should be Findable, Accessible,
Interoperable and Reusable \citep{wilkinson2016fair} -- provide a natural design target for closing
this gap: rather than re-reading primary literature for every comparative question, the evidence
itself should be captured once, in a documented schema, with explicit provenance and with
disagreement preserved rather than silently resolved.

Here we present the \textbf{Fungal BGC Atlas}, a resource that applies this design to 609 fungal
BGC dossiers, each linked to a MIBiG accession. Every dossier was built by large-language-model
(LLM)-assisted extraction from the primary literature and MIBiG record, followed by full manual
review of all 609 dossiers by the curator (Methods). The Atlas resolves each dossier into its
constituent genes, metabolites, experiments, measurements, claims and relationships, all linked back
to their supporting sources, and explicitly records the facets, uncertainties and unresolved
normalisation questions that a single-value schema would otherwise discard. We describe the
resource's construction and content, summarise its composition, and present an interactive web
dashboard and archived data/code package that make the Atlas directly usable for cross-BGC
comparison and genome-mining hypothesis generation.
